\chapter{Conclusion}
\label{ch:conclusion}

\section{Takeaways}

An NLA explanation should be read as two things of different reliability. The gist is what the
reconstruction objective pays for, and it is usually right. A specific name is something the objective
hardly prices at all once the model has moved past the token, so the verbalizer supplies one that fits the
gist. Far from the entity, a name in an explanation is wrong about three times in four.

The verbalizer does, however, know which of its names are shaky. Its own token probabilities separate true
names from wrong ones, and agreement between a few samples does so a little better. Training on that signal
removes some wrong names; training on reconstruction removes none; and no reward tried here brings a lost
specific back.

For someone using an NLA today, the advice is short: trust the gist, and do not trust a name unless several
samples agree on it.

\section{Limitations}
\label{ch:limitations}

\textbf{One model, one layer.} Everything is measured on the released NLA for Qwen2.5-7B at layer~20. The
argument about incentives is general, but the numbers are not, and NLAs for larger models may fade more
slowly.

\textbf{Sixty entities.} The entities are common ones, on Wikipedia text. Rare entities, numbers, dates and
other kinds of specific are untested.

\textbf{The name detector.} Its accuracy has not been measured on a labelled sample, only inspected by hand
on 30 cases. Every comparison uses the same detector, so differences between conditions are safer than
absolute rates.

\textbf{The probe.} It is a floor on what the activation holds, and it chooses among ten known answers while
the verbalizer writes free text. The gap between them is therefore an estimate, not a measurement, of what
the verbalizer leaves unsaid.

\textbf{Small training runs.} Each fine-tuning run uses 2{,}000 activations and a few hundred gradient steps.
A longer run, or true policy-gradient training as the authors used, might move the verbalizer further than
the 8\% seen here.

\section{Next Steps}

The cleanest follow-up is to put the cost on wrong names into the original training, not a small fine-tune
after it: add the agreement penalty to the reward during the authors' reinforcement-learning stage and see
whether the fade in Chapter~\ref{ch:fading} narrows towards the probe. The second is to expose the
confidence the verbalizer already has, by marking low-probability names in the explanation itself.
